\Problem{Suma dels $k$ més petits}


\Statement


Els Jordis (Cortadella i Petit, no els de debò) tenen la conversa següent:

\begin{itemize}

\item[---] Necessitem una funció @int sum_smallest (vector<int> v, int k)@
que, donat un vector $v$ amb $n$ enters i un valor $k\in[0..n]$,
retorni la suma dels $k$ valors més petits de $v$.

\item[---] Vols dir, per exemple, que per $v=\{4,91929,2,4,1\}$ i $k=3$ retorni 7?

\item[---] Exacte!

\item[---] Val. Umhhh... Crec que ja ho tinc: Ordenarem el vector amb un quick sort
i després calcularem la suma de les seves $k$ primeres posicions.

\item[---] És clar! A les transpes d'AP2 i a \textsf{\small lliçons.jutge.org} ja hi tenim una
implementació del  quick sort amb la cèlebre partició de Hoare. Amb un
\emph{cut and paste} ho tindrem solucionat en un momentet.

\item[] \hspace*{-1cm}[10 minuts més tard, el codi complet està llest; vegeu-lo al darrere.]

\item[---] Mira, ja ho tinc fet!

\item[---] Fantàstic, però quin és el cost d'aquest algorisme?

\item[---] Doncs... la fase d'ordenació amb quick sort és $\text{O}(n\log n)$ en mitjana
i la fase de suma és $\text{O}(n)$ perquè $k\le n$. Per tant, el cost total és $\text{O}(n\log n)$ en mitjana.

\item[---] Correcte! Però... cal ordenar tot el vector? No es podria fer més eficient?

\end{itemize}

Modifiqueu el programa dels Jordis per tal que la funció @sum_smallest()@
tingui cost $\text{O}(n)$ en mitjana.




\ObservationNoMain

\medskip
Als jocs de prova del Jutge, els vectors $v$ seràn molt grans i els seus valors seràn
generats aleatòriament.


\clearpage
{\small
\lstinputlisting{public.cc}
}

\clearpage

